% Generated draft: requires mathematical and source review.
\documentclass[11pt]{article}
\usepackage[margin=25mm]{geometry}
\usepackage{fontspec}
\setmainfont{Latin Modern Roman}
\usepackage{amsmath,amssymb,mathtools,unicode-math}
\setmathfont{Latin Modern Math}
\usepackage{enumitem,needspace,xurl,hyperref}
\setlength{\parindent}{0pt}
\setlength{\parskip}{5pt}
\setlength{\emergencystretch}{3em}
\newcommand{\N}{\mathbb{N}}
\newcommand{\Z}{\mathbb{Z}}
\newcommand{\Q}{\mathbb{Q}}
\newcommand{\R}{\mathbb{R}}
\newcommand{\C}{\mathbb{C}}
\newcommand{\F}{\mathbb{F}}
\newcommand{\A}{\mathbb{A}}
\newcommand{\PP}{\mathbb{P}}
\newcommand{\E}{\mathbb{E}}
\newcommand{\eps}{\varepsilon}
\newcommand{\dd}{\,\mathrm{d}}
\begin{document}
\section*{UnsolvedMath ID: 20001900}
\textbf{Source problem code:} catalogue entry
\subsection{Definitions and mathematical statement}
The record merges several page-14 passages: a continuation of Question6.2.1, Section6.3 on flat closing, Section6.4 on entropy rigidity and the start of Section6.5. The original Question6.2.1 is not included. Its continuation asks what follows if all sectional curvatures along a geodesic vanish, or only those of planes containing its tangent, without assuming nonpositive curvature.

Let \(\gamma\) be a unit-speed geodesic in a smooth Riemannian \(n\)-manifold, \(m=n-1\), and let \(e_1(t),\ldots,e_m(t)\) be a parallel orthonormal normal frame. Define the normal Jacobi operator \(A_t(v)=R(v,\dot\gamma)\dot\gamma\), with the usual sign convention. The aid reports equivalence of:
\begin{enumerate}
\item every sectional curvature \(K(\dot\gamma,v)\), \(v\perp\dot\gamma\), vanishes along \(\gamma\);
\item \(A_t=0\) for every \(t\);
\item the normal Jacobi equation in the parallel frame is \(J''=0\), so solutions are parallel transports of \(a+tb\);
\item the space of parallel Jacobi fields along \(\gamma\), including the tangent field, has maximal dimension \(n\).
\end{enumerate}
No curvature-sign hypothesis is used. The finite sectional test consists of the \(m\) directions \(e_i\) and the \(\binom m2\) directions \((e_i+e_j)/\sqrt2\); vanishing on these \(m(m+1)/2\) planes forces the symmetric Jacobi operator to vanish by polarization.

The aid records that this gives neither full transverse sectional flatness nor a flat neighborhood. A product geodesic along the \(\R\) factor of \(\R\times S^2\) has zero radial Jacobi operator despite positive transverse sphere curvature. A warped-torus model \(\dd x^2+f(x)^2\dd y^2\), with \(f'(0)=f''(0)=0\) but \(f''\) nonzero arbitrarily near zero, has a geodesic at \(x=0\) with zero curvature and no flat neighborhood. For example \(f=1+\sin^4x\) supplies these local conditions. These examples use the unrestricted-curvature reading; no nonpositive-curvature counterexample or closing theorem is inferred.

Section6.3 asks Eberlein's flat-closing question: if a compact rank-one nonpositively curved manifold contains a totally geodesic isometric immersion of \(\R^2\), must it contain a totally geodesic isometric immersion of a flat \(T^2\)? Rank one means that some geodesic has rank one. The source reports an affirmative real-analytic result of Bangert--Schroeder and calls the smooth case open. The aid repeats that general smooth flat closing is unresolved in its literature assessment, while mentioning codimension-one and fat-flat positive cases without their precise hypotheses. The Jacobi-operator equivalence does not settle this global closing problem.

Section6.4 asks whether Besson--Courtois--Gallot rigidity extends to higher-rank noncompact-type symmetric spaces. The referenced Theorem5.3.2 is not included. In the standard minimal-entropy interpretation the scale-invariant quantity is \(h_{\mathrm{vol}}(g)^n\operatorname{Vol}(M,g)\), where volume entropy is the exponential universal-cover ball-growth rate. The source credits Connell--Farb's barycenter method for products of negatively curved symmetric factors of dimension at least three, and reports Merlin's 2014 preprint for minimization of \(h_{\mathrm{vol}}^4\operatorname{Vol}\) on compact quotients of \(\mathbb H^2\times\mathbb H^2\). The normalizations and equality/rigidity hypotheses must be taken from those references; arbitrary product-factor scalings are not asserted optimal. The aid regards general irreducible higher-rank minimal-entropy rigidity as still open.

Finally Section6.5 starts with a closed nonpositively curved surface of genus at least two and a count \(N(T)\) of free homotopy classes containing closed geodesics with Gaussian curvature zero along them. The extract stops without giving the dependence on \(T\), a length cutoff or the ensuing question, so no asymptotic counting assertion is reconstructed.
\subsection{Short English statement}
The merged record contains several separate geodesic problems. The aid gives a finite Jacobi-operator test for zero tangent-plane curvature, which does not imply transverse flatness or closing. The flat-closing and higher-rank entropy questions retain their source-specific qualifications, and the final counting fragment is incomplete.
\subsection{Sources}
\begin{itemize}
\item Local UnsolvedMath registry, permanent ID \(20001900\), code catalogue entry; canonical \texttt{aim-geometry-notes.json}, \texttt{notes[237]}; source tag ``question.''
\item AIM geodesics problem list, page14 (PDF index13), continuation of6.2.1 and Sections6.3--6.5, \url{https://aimath.org/pastworkshops/geodesicsproblems.pdf}. The supplied verification identifies the merged boundaries; no PDF inspection is claimed here.
\item Source references: Bangert--Schroeder [Ban-Sch1991]; Besson--Courtois--Gallot [Be-Co-Ga1995] and Theorem5.3.2 in the problem list; Connell--Farb [Con-Far2003a], [Con-Far2003b] and survey [Con-Far2003c]; Merlin [Merlin2014]. Full entries are not supplied and none are invented.
\item Research attempt 1: ``partial\_result,'' ``partially\_solved,'' rigor label ``proved,'' novelty confidence ``low.'' The machine-generated aid is not peer reviewed; independently verify before citation.
\end{itemize}
\end{document}
